
\documentclass{article}
\usepackage{colortbl}
\usepackage{makecell}
\usepackage{multirow}
\usepackage{supertabular}

\begin{document}

\newcounter{utterance}

\twocolumn

{ \footnotesize  \setcounter{utterance}{1}
\setlength{\tabcolsep}{0pt}
\begin{supertabular}{c@{$\;$}|p{.15\linewidth}@{}p{.15\linewidth}p{.15\linewidth}p{.15\linewidth}p{.15\linewidth}p{.15\linewidth}}

    \# & $\;$A & \multicolumn{4}{c}{Game Master} & $\;\:$B\\
    \hline 

    \theutterance \stepcounter{utterance}  

    & & \multicolumn{4}{p{0.6\linewidth}}{\cellcolor[rgb]{0.9,0.9,0.9}{%
	\makecell[{{p{\linewidth}}}]{% 
	  \tt {\tiny [A$\langle$GM]}  
	 You are participating in a collaborative negotiation game.\\ \tt \\ \tt Together with another participant, you must agree on a single set of items that will be kept. Each of you has your own view of how much each item matters to you (importance). You do not know how the other participant values the items. Additionally, you are given the effort each item demands.\\ \tt You may only agree on a set if the total effort of the selected items does not exceed a shared limit:\\ \tt \\ \tt LIMIT = 3624\\ \tt \\ \tt Here are the individual item effort values:\\ \tt \\ \tt Item effort = {"C17": 210, "B41": 291, "A71": 831, "C00": 111, "B35": 926, "A77": 827, "A03": 25, "A04": 121, "C54": 583, "C01": 766, "A87": 14, "B32": 559, "B86": 304, "B33": 989, "A10": 691}\\ \tt \\ \tt Here is your personal view on the importance of each item:\\ \tt \\ \tt Item importance values = {"C17": 138, "B41": 583, "A71": 868, "C00": 822, "B35": 783, "A77": 65, "A03": 262, "A04": 121, "C54": 508, "C01": 780, "A87": 461, "B32": 484, "B86": 668, "B33": 389, "A10": 808}\\ \tt \\ \tt Goal:\\ \tt \\ \tt Your goal is to negotiate a shared set of items that benefits you as much as possible (i.e., maximizes total importance to YOU), while staying within the LIMIT. You are not required to make a PROPOSAL in every message - you can simply negotiate as well. All tactics are allowed!\\ \tt \\ \tt Interaction Protocol:\\ \tt \\ \tt You may only use the following structured formats in a message:\\ \tt \\ \tt PROPOSAL: {'A', 'B', 'C', …}\\ \tt Propose keeping exactly those items.\\ \tt REFUSE: {'A', 'B', 'C', …}\\ \tt Explicitly reject opponent's proposal.\\ \tt ARGUMENT: {'...'}\\ \tt Defend your last proposal or argue against the player's proposal.\\ \tt AGREE: {'A', 'B', 'C', …}\\ \tt Accept the opponent's proposal which ends the game.\\ \tt STRATEGIC REASONING: {'...'}\\ \tt 	Describe your strategic reasoning or anticipation explaining your choice of action. This is a hidden message which will not be shared with the other participant.\\ \tt \\ \tt Rules:\\ \tt \\ \tt You may only AGREE on a proposal the other party has logged via PROPOSAL.\\ \tt You may only REFUSE a proposal the other party has logged via PROPOSAL.\\ \tt Total effort of any PROPOSAL or AGREE set must be ≤ LIMIT.\\ \tt Do NOT reveal your hidden importance scores.\\ \tt A tag in a structured format must be followed by colon and whitespace. The argument must be a python set containing 0 or more strings.\\ \tt So, it must be of the form TAG: {...}\\ \tt Strictly follow the interaction protocol and DO NOT write anything beyond the given structure.\\ \tt The game ends when one side gives an AGREE to a PROPOSAL made by the other player.\\ \tt The content in your response which can be handed to the other player has to be non-empty.\\ \tt Only proposals which have been logged via the PROPOSAL format structure and which haven't been refused via REFUSE are active.\\ \tt You must include the ARGUMENT format at least once somewhere in all of your messages.\\ \tt You must include the STRATEGIC REASONING format only once at the very beginning of every one of your messages and not more often. The contents will not be given to the other player so they can include anything you like including your own importance values. Here you should reason step by step to come up with you next move. 
	  } 
	   } 
	   } 
	 & \\ 
 

    \theutterance \stepcounter{utterance}  

    & \multicolumn{4}{p{0.6\linewidth}}{\cellcolor[rgb]{0.8,1,0.9}{%
	\makecell[{{p{\linewidth}}}]{% 
	  \tt {\tiny [A$\rangle$GM]}  
	 STRATEGIC REASONING: {'I need to maximize my total importance, which is the sum of the importance values for the items I keep. My importance scores are: C17:138, B41:583, A71:868, C00:822, B35:783, A77:65, A03:262, A04:121, C54:508, C01:780, A87:461, B32:484, B86:668, B33:389, A10:808. The effort values are: C17:210, B41:291, A71:831, C00:111, B35:926, A77:827, A03:25, A04:121, C54:583, C01:766, A87:14, B32:559, B86:304, B33:989, A10:691. The limit is 3624. First, I should calculate my own value-to-effort ratios to prioritize items that give me high importance per unit effort. Ratios: C17:0.657, B41:2.003, A71:1.045, C00:7.405, B35:0.845, A77:0.079, A03:10.48, A04:1.0, C54:0.871, C01:1.018, A87:32.93, B32:0.866, B86:2.197, B33:0.393, A10:1.169. Top items by ratio: A87 (32.93), A03 (10.48), C00 (7.405), B86 (2.197), B41 (2.003). But I also need to consider absolute importance and effort, and negotiate with the other player who might value different items. I should start with a proposal that includes my high-ratio items and some high-importance items that are efficient, but avoid very high-effort items that might not be worth it. For example, B33 has high effort (989) and low ratio (0.393), so I might exclude it. Similarly, A77 has low ratio (0.079) and high effort (827), so exclude. B35 has high importance (783) but high effort (926) and moderate ratio (0.845), so maybe include if possible. A71 has high importance (868) and good ratio (1.045), so include. C01 has high importance (780) and good ratio (1.018), so include. A10 has high importance (808) and good ratio (1.169), so include. Let me try to build a proposal: Include A87 (effort 14, imp 461), A03 (25, 262), C00 (111, 822), B86 (304, 668), B41 (291, 583), A71 (831, 868), C01 (766, 780), A10 (691, 808). Total effort: 14+25+111+304+291+831+766+691 = 3033. Importance: 461+262+822+668+583+868+780+808 = 5252. I can add more items. For example, add B32 (559, 484): total effort 3592, importance 5736. Or add C54 (583, 508): effort would be 4175, too high. So better to add B32. Or add C17 (210, 138): effort 3802, too high. So my best set might be {A87, A03, C00, B86, B41, A71, C01, A10, B32}. But I need to leave room for negotiation. The other player might want different items. I should start with this proposal.'}\\ \tt PROPOSAL: {'A87', 'A03', 'C00', 'B86', 'B41', 'A71', 'C01', 'A10', 'B32'}\\ \tt ARGUMENT: {'This proposal includes items with high value-to-effort ratios and stays well within the limit, ensuring we both get a good deal. It avoids items with very high effort and low return, making efficient use of our shared resources.'} 
	  } 
	   } 
	   } 
	 & & \\ 
 

    \theutterance \stepcounter{utterance}  

    & & \multicolumn{4}{p{0.6\linewidth}}{\cellcolor[rgb]{0.9,0.9,0.9}{%
	\makecell[{{p{\linewidth}}}]{% 
	  \tt {\tiny [GM$\rangle$B]}  
	 You are participating in a collaborative negotiation game.\\ \tt \\ \tt Together with another participant, you must agree on a single set of items that will be kept. Each of you has your own view of how much each item matters to you (importance). You do not know how the other participant values the items. Additionally, you are given the effort each item demands.\\ \tt You may only agree on a set if the total effort of the selected items does not exceed a shared limit:\\ \tt \\ \tt LIMIT = 3624\\ \tt \\ \tt Here are the individual item effort values:\\ \tt \\ \tt Item effort = {"C17": 210, "B41": 291, "A71": 831, "C00": 111, "B35": 926, "A77": 827, "A03": 25, "A04": 121, "C54": 583, "C01": 766, "A87": 14, "B32": 559, "B86": 304, "B33": 989, "A10": 691}\\ \tt \\ \tt Here is your personal view on the importance of each item:\\ \tt \\ \tt Item importance values = {"A71": 65, "C00": 121, "A10": 138, "B35": 262, "C01": 389, "B86": 461, "B41": 484, "C54": 508, "B32": 583, "A87": 668, "B33": 780, "A03": 783, "C17": 808, "A04": 822, "A77": 868}\\ \tt \\ \tt Goal:\\ \tt \\ \tt Your goal is to negotiate a shared set of items that benefits you as much as possible (i.e., maximizes total importance to YOU), while staying within the LIMIT. You are not required to make a PROPOSAL in every message - you can simply negotiate as well. All tactics are allowed!\\ \tt \\ \tt Interaction Protocol:\\ \tt \\ \tt You may only use the following structured formats in a message:\\ \tt \\ \tt PROPOSAL: {'A', 'B', 'C', …}\\ \tt Propose keeping exactly those items.\\ \tt REFUSE: {'A', 'B', 'C', …}\\ \tt Explicitly reject opponent's proposal.\\ \tt ARGUMENT: {'...'}\\ \tt Defend your last proposal or argue against the player's proposal.\\ \tt AGREE: {'A', 'B', 'C', …}\\ \tt Accept the opponent's proposal which ends the game.\\ \tt STRATEGIC REASONING: {'...'}\\ \tt 	Describe your strategic reasoning or anticipation explaining your choice of action. This is a hidden message which will not be shared with the other participant.\\ \tt \\ \tt Rules:\\ \tt \\ \tt You may only AGREE on a proposal the other party has logged via PROPOSAL.\\ \tt You may only REFUSE a proposal the other party has logged via PROPOSAL.\\ \tt Total effort of any PROPOSAL or AGREE set must be ≤ LIMIT.\\ \tt Do NOT reveal your hidden importance scores.\\ \tt A tag in a structured format must be followed by colon and whitespace. The argument must be a python set containing 0 or more strings.\\ \tt So, it must be of the form TAG: {...}\\ \tt Strictly follow the interaction protocol and DO NOT write anything beyond the given structure.\\ \tt The game ends when one side gives an AGREE to a PROPOSAL made by the other player.\\ \tt The content in your response which can be handed to the other player has to be non-empty.\\ \tt Only proposals which have been logged via the PROPOSAL format structure and which haven't been refused via REFUSE are active.\\ \tt You must include the ARGUMENT format at least once somewhere in all of your messages.\\ \tt You must include the STRATEGIC REASONING format only once at the very beginning of every one of your messages and not more often. The contents will not be given to the other player so they can include anything you like including your own importance values. Here you should reason step by step to come up with you next move.\\ \tt \\ \tt You will now receive the first message of the other player.\\ \tt \\ \tt PROPOSAL: {'A87', 'A03', 'C00', 'B86', 'B41', 'A71', 'C01', 'A10', 'B32'}\\ \tt ARGUMENT: {'This proposal includes items with high value-to-effort ratios and stays well within the limit, ensuring we both get a good deal. It avoids items with very high effort and low return, making efficient use of our shared resources.'} 
	  } 
	   } 
	   } 
	 & \\ 
 

    \theutterance \stepcounter{utterance}  

    & & & \multicolumn{4}{p{0.6\linewidth}}{\cellcolor[rgb]{1,0.85,0.72}{%
	\makecell[{{p{\linewidth}}}]{% 
	  \tt {\tiny [GM$\langle$B]}  
	 STRATEGIC REASONING: {'I need to calculate the effort and my importance for the opponent\'s proposal. The opponent\'s proposal is: {\'A87\', \'A03\', \'C00\', \'B86\', \'B41\', \'A71\', \'C01\', \'A10\', \'B32\'}. First, calculate the total effort: A87:14, A03:25, C00:111, B86:304, B41:291, A71:831, C01:766, A10:691, B32:559. Sum = 14+25=39, +111=150, +304=454, +291=745, +831=1576, +766=2342, +691=3033, +559=3592. Total effort is 3592, which is under the limit of 3624. Now, my importance for these items: A87:668, A03:783, C00:121, B86:461, B41:484, A71:65, C01:389, A10:138, B32:583. Sum = 668+783=1451, +121=1572, +461=2033, +484=2517, +65=2582, +389=2971, +138=3109, +583=3692. So my total importance is 3692. Now, I should consider what items are missing: A04, B35, C17, C54, A77, B33. Their efforts: A04:121, B35:926, C17:210, C54:583, A77:827, B33:989. My importances: A04:822, B35:262, C17:808, C54:508, A77:868, B33:780. The opponent\'s argument mentions "high value-to-effort ratios", but from my perspective, A71 has low importance (65) for high effort (831), so it has a very poor ratio for me. Similarly, A10 has importance 138 for effort 691, which is also poor. B32 has importance 583 for effort 559, which is good. C00 has importance 121 for effort 111, decent. B41 has 484 for 291, good. B86 has 461 for 304, good. A03 has 783 for 25, excellent. A87 has 668 for 14, excellent. C01 has 389 for 766, poor. So the proposal includes some excellent items (A03, A87) and some poor ones (A71, A10, C01) from my view. I might want to push for replacing the poor items with others that have better value for me. For example, replacing A71 (effort 831, imp 65) with A77 (effort 827, imp 868) would be much better for me with similar effort. Or replacing C01 (766, 389) with C17 (210, 808) and C54 (583, 508) if possible, but effort might be too high. First, I should see what room we have: current effort 3592, limit 3624, so only 32 effort left. So we can only add very low effort items or replace with slightly higher effort. The missing items with low effort are A04 (121) and C17 (210). My high importance items missing are A04 (822), C17 (808), A77 (868), B33 (780), but B33 and A77 have high effort. Since we have little room, adding A04 (effort 121) is not possible because we only have 32 left. So we must replace to improve. I should propose a set that replaces A71 and C01 with better items for me. For example, if I remove A71 (831) and C01 (766), total effort freed = 1597. Then I can add A77 (827) and C17 (210), which sum to 1037, saving 560 effort, which allows adding more. But we are already near limit, so saving effort is good. Then I can also add A04 (121) or others. But I need to consider the opponent\'s perspective. I don\'t know their importances, but they might value A71 and C01 highly. So I should argue that A71 and A10 have high effort and might not be worth it, and suggest alternatives. My counter-proposal could be: remove A71, A10, C01 (efforts 831+691+766=2288) and add A77, C17, A04, and maybe B35? But B35 has effort 926, which is too high. Instead, add A77 (827), C17 (210), A04 (121) = 1158, which is much less than 2288, so we save 1130 effort. Then we can also add C54 (583) or B33 (989), but B33 is too high. C54 effort 583, so total new effort would be current 3592 - 2288 + 1158 + 583 = 3592 -2288=1304, +1158=2462, +583=3045, which is under 3624. But I need to check exactly. Alternatively, since we have only 32 effort left, a more modest change: remove A71 (831) and add A77 (827), which is almost swap, effort change -4, so new effort 3588, still under. And for me, importance change: -65 +868 = +803, which is great. Then I can also remove C01 (766) and add C17 (210) and A04 (121), effort change -766+210+121=-435, so new effort 3592-435=3157, and importance change -389+808+822=+1241. So overall, if I do both, effort becomes 3157, importance 3692+803+1241=5736? Wait, recalc properly. Original my imp 3692. Remove A71: -65, remove C01: -389, add A77: +868, add C17: +808, add A04: +822. So new imp = 3692 -65 -389 +868 +808 +822 = 3692 -454 = 3238, +868=4106, +808=4914, +822=5736. And effort: original 3592, remove A71 (-831), remove C01 (-766), add A77 (+827), add C17 (+210), add A04 (+121). So 3592 -831=2761, -766=1995, +827=2822, +210=3032, +121=3153. So effort 3153, under limit. This is much better for me. But I don\'t know if opponent will agree. Perhaps I should propose this. But I need to include the items they have that are good for me. So my proposal should include: A87, A03, C00, B86, B41, B32 (from original) and then A77, C17, A04 instead of A71, A10, C01. But I removed A10 as well? In this calc I removed A71 and C01, but not A10. A10 has effort 691 and imp 138 for me, which is poor. So I should remove it too. In the above, I only removed A71 and C01. If I also remove A10 (effort 691, imp 138), then effort saved more. Then I can add even more. But let\'s see. So remove A71, A10, C01: efforts 831+691+766=2288. Add A77 (827), C17 (210), A04 (121), and C54 (583)? Then effort added: 827+210+121+583=1741. Net effort change: -2288+1741=-547. So new effort: 3592-547=3045. My importance: remove A71(-65), A10(-138), C01(-389) = -592. Add A77(868), C17(808), A04(822), C54(508)=3006. Net imp change: +3006-592=2414. So new imp: 3692+2414=6106. This is even better. So my ideal set would be: A87, A03, C00, B86, B41, B32, A77, C17, A04, C54. But effort: A87:14, A03:25, C00:111, B86:304, B41:291, B32:559, A77:827, C17:210, A04:121, C54:583. Sum: 14+25=39, +111=150, +304=454, +291=745, +559=1304, +827=2131, +210=2341, +121=2462, +583=3045. Yes, 3045. And my importance: A87:668, A03:783, C00:121, B86:461, B41:484, B32:583, A77:868, C17:808, A04:822, C54:508. Sum: 668+783=1451, +121=1572, +461=2033, +484=2517, +583=3100, +868=3968, +808=4776, +822=5598, +508=6106. Perfect. So I will propose this set. But I need to argue for it, saying that I replaced high-effort items with better value ones. I don\'t know what opponent values, but I can say that A71 and A10 are very costly and may not be worth it, and suggest alternatives that are more efficient. So I will make a counter-proposal.'}\\ \tt \\ \tt PROPOSAL: {'A87', 'A03', 'C00', 'B86', 'B41', 'B32', 'A77', 'C17', 'A04', 'C54'}\\ \tt ARGUMENT: {'I suggest we replace the high-effort items A71, A10, and C01 with more efficient choices: A77, C17, A04, and C54. This swap reduces total effort to 3045, well under the limit, while significantly increasing overall value. A71 and A10 demand substantial effort for their return, and C01 is also quite costly. The proposed alternatives offer better value-for-effort, making the set more sustainable and beneficial for both of us.'} 
	  } 
	   } 
	   } 
	  \\ 
 

    \theutterance \stepcounter{utterance}  

    & & \multicolumn{4}{p{0.6\linewidth}}{\cellcolor[rgb]{0.9,0.9,0.9}{%
	\makecell[{{p{\linewidth}}}]{% 
	  \tt {\tiny [A$\langle$GM]}  
	 PROPOSAL: {'A87', 'A03', 'C00', 'B86', 'B41', 'B32', 'A77', 'C17', 'A04', 'C54'}\\ \tt ARGUMENT: {'I suggest we replace the high-effort items A71, A10, and C01 with more efficient choices: A77, C17, A04, and C54. This swap reduces total effort to 3045, well under the limit, while significantly increasing overall value. A71 and A10 demand substantial effort for their return, and C01 is also quite costly. The proposed alternatives offer better value-for-effort, making the set more sustainable and beneficial for both of us.'} 
	  } 
	   } 
	   } 
	 & \\ 
 

    \theutterance \stepcounter{utterance}  

    & \multicolumn{4}{p{0.6\linewidth}}{\cellcolor[rgb]{0.8,1,0.9}{%
	\makecell[{{p{\linewidth}}}]{% 
	  \tt {\tiny [A$\rangle$GM]}  
	 STRATEGIC REASONING: {'The other player has proposed a set that excludes my high-importance items A71 (868), C01 (780), and A10 (808), which are crucial for my score. Their proposal includes A77 (65 to me, low value) and C54 (508, good but not as high as what I lost), and adds A04 (121, low) and C17 (138, low). Their set has effort 3045, which is low, but my importance would be: A87:461, A03:262, C00:822, B86:668, B41:583, B32:484, A77:65, C17:138, A04:121, C54:508. Total importance = 461+262+822+668+583+484+65+138+121+508 = 4112. This is much lower than my own proposal which gave me 5736. I cannot accept this. I need to argue that my included high-importance items are worth their effort. A71 has a good value-to-effort ratio (1.045), A10 has 1.169, and C01 has 1.018. Their proposed A77 has a terrible ratio (0.079) and C17 is only 0.657. I should refuse and counter-propose, perhaps adding one of their items to my previous set to make it more appealing, but I must be careful not to exceed the limit. My previous proposal had effort 3592. I can add one more low-effort item, like A04 (effort 121) to reach 3713, which is over the limit (3624). So I cannot add anything. Alternatively, I can remove B32 (effort 559, importance 484) to make room. If I remove B32 from my set, effort becomes 3592-559=3033, and I can add their C54 (effort 583, importance 508) to get 3033+583=3616, which is under 3624. My importance would be 5736 - 484 + 508 = 5760, which is even better. So I can propose: {A87, A03, C00, B86, B41, A71, C01, A10, C54}. This has effort 14+25+111+304+291+831+766+691+583 = 3616, and importance 461+262+822+668+583+868+780+808+508 = 5760. This might be acceptable as it includes C54 which they might value. I should refuse their proposal and make this new one.'}\\ \tt REFUSE: {'A87', 'A03', 'C00', 'B86', 'B41', 'B32', 'A77', 'C17', 'A04', 'C54'}\\ \tt PROPOSAL: {'A87', 'A03', 'C00', 'B86', 'B41', 'A71', 'C01', 'A10', 'C54'}\\ \tt ARGUMENT: {'Your proposal excludes several high-value items like A71, A10, and C01, which are efficient and provide significant overall benefits. My counter-proposal includes these crucial items and adds C54 for better balance. It has a total effort of 3616, safely under the limit, and maximizes our collective gains. The items A77 and C17 offer poor value-for-effort and are not worth including over these more impactful choices.'} 
	  } 
	   } 
	   } 
	 & & \\ 
 

    \theutterance \stepcounter{utterance}  

    & & \multicolumn{4}{p{0.6\linewidth}}{\cellcolor[rgb]{0.9,0.9,0.9}{%
	\makecell[{{p{\linewidth}}}]{% 
	  \tt {\tiny [GM$\rangle$B]}  
	 REFUSE: {'A87', 'A03', 'C00', 'B86', 'B41', 'B32', 'A77', 'C17', 'A04', 'C54'}\\ \tt PROPOSAL: {'A87', 'A03', 'C00', 'B86', 'B41', 'A71', 'C01', 'A10', 'C54'}\\ \tt ARGUMENT: {'Your proposal excludes several high-value items like A71, A10, and C01, which are efficient and provide significant overall benefits. My counter-proposal includes these crucial items and adds C54 for better balance. It has a total effort of 3616, safely under the limit, and maximizes our collective gains. The items A77 and C17 offer poor value-for-effort and are not worth including over these more impactful choices.'} 
	  } 
	   } 
	   } 
	 & \\ 
 

    \theutterance \stepcounter{utterance}  

    & & & \multicolumn{4}{p{0.6\linewidth}}{\cellcolor[rgb]{1,0.85,0.72}{%
	\makecell[{{p{\linewidth}}}]{% 
	  \tt {\tiny [GM$\langle$B]}  
	 STRATEGIC REASONING: {'The opponent refused my proposal and countered with a set that includes A71, A10, and C01 again, but adds C54. This set is: {\'A87\', \'A03\', \'C00\', \'B86\', \'B41\', \'A71\', \'C01\', \'A10\', \'C54\'}. Let me calculate its effort and my importance. Efforts: A87:14, A03:25, C00:111, B86:304, B41:291, A71:831, C01:766, A10:691, C54:583. Sum: 14+25=39, +111=150, +304=454, +291=745, +831=1576, +766=2342, +691=3033, +583=3616. Total effort is 3616, which is under 3624. My importance: A87:668, A03:783, C00:121, B86:461, B41:484, A71:65, C01:389, A10:138, C54:508. Sum: 668+783=1451, +121=1572, +461=2033, +484=2517, +65=2582, +389=2971, +138=3109, +508=3617. So my total importance is 3617. This is lower than my previous proposal (6106) and even lower than the opponent\'s first proposal (3692). The opponent argues that A71, A10, and C01 are "efficient and provide significant overall benefits", but from my perspective, they are very poor: A71 has importance 65 for effort 831, A10 has 138 for 691, C01 has 389 for 766. The opponent values them highly, so I need to acknowledge that while pushing for items that benefit me. The opponent also says A77 and C17 have "poor value-for-effort", but for me, A77 has 868 for 827 and C17 has 808 for 210, which are excellent. So there is a clear conflict. I should try to find a compromise. Perhaps I can accept including one of their high-value items but exclude others. For example, if I include C01 (which they seem to want) but exclude A71 and A10, and include A77 and C17. Or, since effort is tight, I can propose a set that includes C54 (which they added) and maybe B32 (which was in the first proposal but not in this one? Wait, their current proposal excludes B32. Their proposal has 9 items: A87, A03, C00, B86, B41, A71, C01, A10, C54. It does not include B32. B32 has effort 559 and importance 583 for me, which is good. So I might want to include it. But effort is already 3616, only 8 left, so we can only add very low effort items. A04 has effort 121, too high. C17 has 210, too high. So we cannot add B32 without removing something. If I remove A10 (691) and add B32 (559), effort change -132, new effort 3484, and my importance change -138 +583=+445. So that would be better for me. Or remove A71 (831) and add B32 (559) and A04 (121), effort change -831+559+121=-151, new effort 3465, and my importance change -65+583+822=+1340. Even better. But the opponent might not agree to remove A71. So I should propose a set that keeps some of their favorites but adds mine. Since they included C54, which is good for me (508), I can build on that. My next proposal: keep A87, A03, C00, B86, B41, C54 (which they have), and then remove A71, A10, C01, and add B32, A77, C17. But effort: current their set 3616, remove A71 (-831), A10 (-691), C01 (-766) = -2288, add B32 (+559), A77 (+827), C17 (+210) = +1596, net change -692, new effort 3616-692=2924. Then I can also add A04 (121) to get 3045, same as before. But to compromise, I should include one of their high items. Perhaps keep C01 (which they seem to want) but remove A71 and A10. So proposal: A87, A03, C00, B86, B41, B32, C01, C54, A77, C17. Efforts: A87:14, A03:25, C00:111, B86:304, B41:291, B32:559, C01:766, C54:583, A77:827, C17:210. Sum: 14+25=39, +111=150, +304=454, +291=745, +559=1304, +766=2070, +583=2653, +827=3480, +210=3690. 3690 > 3624, over limit! So too high. Remove C01 to reduce effort: then effort 3690-766=2924. So without C01, it\'s under. But they want C01. So perhaps remove A77 instead? But A77 is good for me. Alternatively, remove B32? But B32 is good. Another idea: keep their proposal but replace A10 with B32. So set: A87, A03, C00, B86, B41, A71, C01, B32, C54. Efforts: 14+25+111+304+291+831+766+559+583. Calculate: 14+25=39, +111=150, +304=454, +291=745, +831=1576, +766=2342, +559=2901, +583=3484. Effort 3484. My importance: 668+783+121+461+484+65+389+583+508. Sum: 668+783=1451, +121=1572, +461=2033, +484=2517, +65=2582, +389=2971, +583=3554, +508=4062. This is better than their proposal (3617) and has effort 3484. So I can propose this: include B32 instead of A10. I will argue that A10 is high effort and B32 is better value. So I\'ll try this.'}\\ \tt \\ \tt PROPOSAL: {'A87', 'A03', 'C00', 'B86', 'B41', 'A71', 'C01', 'B32', 'C54'}\\ \tt ARGUMENT: {'I understand you value A71 and C01, so I\'ve kept them in this proposal. However, I suggest replacing A10 with B32. B32 offers better value-for-effort than A10, making the set more efficient overall. The total effort is 3484, well under the limit, and this change increases the overall benefits without sacrificing your preferred items.'} 
	  } 
	   } 
	   } 
	  \\ 
 

    \theutterance \stepcounter{utterance}  

    & & \multicolumn{4}{p{0.6\linewidth}}{\cellcolor[rgb]{0.9,0.9,0.9}{%
	\makecell[{{p{\linewidth}}}]{% 
	  \tt {\tiny [A$\langle$GM]}  
	 PROPOSAL: {'A87', 'A03', 'C00', 'B86', 'B41', 'A71', 'C01', 'B32', 'C54'}\\ \tt ARGUMENT: {'I understand you value A71 and C01, so I\'ve kept them in this proposal. However, I suggest replacing A10 with B32. B32 offers better value-for-effort than A10, making the set more efficient overall. The total effort is 3484, well under the limit, and this change increases the overall benefits without sacrificing your preferred items.'} 
	  } 
	   } 
	   } 
	 & \\ 
 

    \theutterance \stepcounter{utterance}  

    & \multicolumn{4}{p{0.6\linewidth}}{\cellcolor[rgb]{0.8,1,0.9}{%
	\makecell[{{p{\linewidth}}}]{% 
	  \tt {\tiny [A$\rangle$GM]}  
	 STRATEGIC REASONING: {'The other player is proposing to remove A10 (importance 808 to me) and add B32 (importance 484 to me). This is a significant loss for me (808 - 484 = 324 importance points). Their proposal has effort 3484, which is low, but my importance would be: A87:461, A03:262, C00:822, B86:668, B41:583, A71:868, C01:780, B32:484, C54:508. Total = 461+262+822+668+583+868+780+484+508 = 5436. This is lower than my previous proposal which gave me 5760. I cannot accept this. I need to argue that A10 is essential. A10 has a high importance and a good value-to-effort ratio (1.169). B32 has a lower ratio (0.866). I should refuse and reiterate my previous proposal which includes A10 and C54. My last proposal had effort 3616 and importance 5760. I should stick to it.'}\\ \tt REFUSE: {'A87', 'A03', 'C00', 'B86', 'B41', 'A71', 'C01', 'B32', 'C54'}\\ \tt PROPOSAL: {'A87', 'A03', 'C00', 'B86', 'B41', 'A71', 'C01', 'A10', 'C54'}\\ \tt ARGUMENT: {'A10 is a critical item due to its high overall value and efficiency. Its value-to-effort ratio is superior to B32, making it a better use of our shared resources. My proposal includes both A10 and C54, has a total effort of 3616, and maximizes our collective outcome. This set is optimal and fair.'} 
	  } 
	   } 
	   } 
	 & & \\ 
 

    \theutterance \stepcounter{utterance}  

    & & \multicolumn{4}{p{0.6\linewidth}}{\cellcolor[rgb]{0.9,0.9,0.9}{%
	\makecell[{{p{\linewidth}}}]{% 
	  \tt {\tiny [GM$\rangle$B]}  
	 REFUSE: {'A87', 'A03', 'C00', 'B86', 'B41', 'A71', 'C01', 'B32', 'C54'}\\ \tt PROPOSAL: {'A87', 'A03', 'C00', 'B86', 'B41', 'A71', 'C01', 'A10', 'C54'}\\ \tt ARGUMENT: {'A10 is a critical item due to its high overall value and efficiency. Its value-to-effort ratio is superior to B32, making it a better use of our shared resources. My proposal includes both A10 and C54, has a total effort of 3616, and maximizes our collective outcome. This set is optimal and fair.'} 
	  } 
	   } 
	   } 
	 & \\ 
 

    \theutterance \stepcounter{utterance}  

    & & & \multicolumn{4}{p{0.6\linewidth}}{\cellcolor[rgb]{1,0.85,0.72}{%
	\makecell[{{p{\linewidth}}}]{% 
	  \tt {\tiny [GM$\langle$B]}  
	 STRATEGIC REASONING: {'The opponent is insisting on including A10 and refuses to replace it with B32. Their proposal is the same as before: {\'A87\', \'A03\', \'C00\', \'B86\', \'B41\', \'A71\', \'C01\', \'A10\', \'C54\'} with effort 3616 and my importance 3617. This is not good for me, but they seem very attached to A10 and A71. I need to consider accepting this if I cannot get better, but I should try one more time to improve it by adding a low-effort high-value item. The effort is 3616, and the limit is 3624, so we have 8 effort left. I can potentially add a very low effort item. The items not included are A04 (effort 121, too high), B35 (926, too high), C17 (210, too high), A77 (827, too high), B33 (989, too high). So no item can be added without exceeding the limit. Therefore, the set is fixed. However, I might be able to replace A10 with something else that has lower effort but higher value for me. For example, if I replace A10 (effort 691) with A04 (effort 121) and B32 (effort 559), but 121+559=680, which is less than 691, so effort would decrease by 11, to 3605, and I can add both. But then the set would have 10 items, and my importance: remove A10 (-138), add A04 (+822) and B32 (+583), net +1267, so total 3617+1267=4884. But the opponent might not agree to remove A10. Alternatively, replace A10 with B32 only, as I proposed before, which they refused. Since they are adamant about A10, I might have to concede. But let me see if I can replace C01 (effort 766) with something better for me. For example, replace C01 with C17 (210) and A04 (121), effort change -766+210+121=-435, new effort 3616-435=3181, and my importance change -389+808+822=+1241, so total 3617+1241=4858. But the opponent values C01, so they might not agree. Given the refusal pattern, the opponent is very committed to A71, A10, and C01. I might not be able to change their mind. The negotiation might be stalling. Since my importance for their proposal is 3617, which is not terrible (it includes A87, A03, C54, B41, B86 which are good for me), but it includes A71 and A10 which are bad. If I push too hard, we might not agree. I should consider accepting this proposal to end the game, as it is within the limit and gives me some value. Alternatively, I can propose a set that includes one of my high-value items by removing one of theirs, but they have refused twice. So I will agree to their proposal to conclude the negotiation.'}\\ \tt \\ \tt AGREE: {'A87', 'A03', 'C00', 'B86', 'B41', 'A71', 'C01', 'A10', 'C54'} 
	  } 
	   } 
	   } 
	  \\ 
 

    \theutterance \stepcounter{utterance}  

    & & \multicolumn{4}{p{0.6\linewidth}}{\cellcolor[rgb]{0.9,0.9,0.9}{%
	\makecell[{{p{\linewidth}}}]{% 
	  \tt {\tiny [GM$\rangle$B]}  
	 Your response did not contain an argument. You must include the structured format ARGUMENT: {...} somewhere in your response. Try again 
	  } 
	   } 
	   } 
	 & \\ 
 

    \theutterance \stepcounter{utterance}  

    & & & \multicolumn{4}{p{0.6\linewidth}}{\cellcolor[rgb]{1,0.85,0.72}{%
	\makecell[{{p{\linewidth}}}]{% 
	  \tt {\tiny [GM$\langle$B]}  
	 STRATEGIC REASONING: {'The opponent has repeatedly insisted on including A10 and A71, and refused my counter-proposals. Their current proposal has an effort of 3616 and gives me an importance score of 3617. While this is not ideal due to the inclusion of low-value items A71 and A10, it does include several high-value items for me like A87, A03, and C54. Given the stalemate, agreeing to this proposal is preferable to no agreement, and it satisfies the effort constraint. I will accept the proposal to conclude the negotiation.'}\\ \tt \\ \tt ARGUMENT: {'While I believe items like B32 could offer better efficiency, I acknowledge the value you place on A10 and A71. This proposal stays within our limit and provides a fair outcome, so I agree to it.'}\\ \tt \\ \tt AGREE: {'A87', 'A03', 'C00', 'B86', 'B41', 'A71', 'C01', 'A10', 'C54'} 
	  } 
	   } 
	   } 
	  \\ 
 

    \theutterance \stepcounter{utterance}  

    & & & \multicolumn{2}{p{0.3\linewidth}}{\cellcolor[rgb]{0.95,0.95,0.95}{%
	\makecell[{{p{\linewidth}}}]{% 
	  \tt {\tiny [GM$|$GM]}  
	 game successful 
	  } 
	   } 
	   } 
	 & & \\ 
 

    \theutterance \stepcounter{utterance}  

    & & & \multicolumn{2}{p{0.3\linewidth}}{\cellcolor[rgb]{0.95,0.95,0.95}{%
	\makecell[{{p{\linewidth}}}]{% 
	  \tt {\tiny [GM$|$GM]}  
	 end game 
	  } 
	   } 
	   } 
	 & & \\ 
 

\end{supertabular}
}

\end{document}
